\chapter{Operators}

\section{Arithmetic}

\begin{center}
\begin{tabular}{cll}
\toprule
\textbf{Operator} & \textbf{Operation} & \textbf{Example} \\
\midrule
\texttt{+}  & Addition       & \texttt{3 + 2} → \texttt{5} \\
\texttt{-}  & Subtraction    & \texttt{5 - 3} → \texttt{2} \\
\texttt{*}  & Multiplication & \texttt{4 * 3} → \texttt{12} \\
\texttt{/}  & Division       & \texttt{7 / 2} → \texttt{3} (integers) \\
\texttt{\^} & Power          & \texttt{2\^{}10} → \texttt{1024} \\
\texttt{\%} & Modulo         & \texttt{7 \% 3} → \texttt{1} \\
\bottomrule
\end{tabular}
\end{center}

Division between two \texttt{Int} values is integer division (truncated
toward zero). For real division, at least one operand must be
\texttt{Float}:

\begin{lstlisting}
7 / 2      // 3
7 / 2.0    // 3.5
7.0 / 2    // 3.5
\end{lstlisting}

\section{Comparison}

\begin{center}
\begin{tabular}{cll}
\toprule
\textbf{Operator} & \textbf{Meaning} & \textbf{Example} \\
\midrule
\texttt{==} & Equal to              & \texttt{3 == 3} → \texttt{true} \\
\texttt{!=} & Not equal to          & \texttt{3 != 4} → \texttt{true} \\
\texttt{>}  & Greater than          & \texttt{5 > 3}  → \texttt{true} \\
\texttt{<}  & Less than             & \texttt{3 < 5}  → \texttt{true} \\
\texttt{>=} & Greater than or equal & \texttt{3 >= 3} → \texttt{true} \\
\texttt{<=} & Less than or equal    & \texttt{2 <= 3} → \texttt{true} \\
\bottomrule
\end{tabular}
\end{center}

\section{Logical}

\begin{lstlisting}
true and false    // false
true or false     // true
not true          // false
\end{lstlisting}

\lstinline{and} and \lstinline{or} use short-circuit evaluation: the
second operand is not evaluated if the result is already determined by
the first.

\section{The \texttt{in} operator}

Tests membership in a list:

\begin{lstlisting}
let regions = ["north", "south", "center"]
"south" in regions    // true
"east"  in regions    // false
\end{lstlisting}

Also works with ranges:

\begin{lstlisting}
5  in 1..10    // true
15 in 1..10    // false
\end{lstlisting}

\section{The pipe operator \texttt{|>}}
\label{sec:pipe}

The \lstinline{|>} operator passes the value on the left as the first
argument to the function on the right:

\begin{lstlisting}
// without pipe
display filter(df, x > 5)

// with pipe — more readable
df |> filter(x > 5) |> display
\end{lstlisting}

The pipe has special semantics with DataFrames: when used without an
explicit assignment (\lstinline{let}), the result is assigned back to the
source variable:

\begin{lstlisting}
let df = load("data.csv")

// updates df in place
df |> filter(x > 0) |> sort(x)

// creates df2, df remains unchanged
let df2 = df |> filter(x > 0)
\end{lstlisting}

This behaviour is discussed in detail in Chapter~\ref{cap:pipes}.

\section{Precedence}

From lowest to highest priority:

\begin{enumerate}
  \item \texttt{or}
  \item \texttt{and}
  \item \texttt{not}
  \item \texttt{== != > < >= <=}
  \item \texttt{+ -}
  \item \texttt{* / \%}
  \item \texttt{\^} (right-associative)
  \item Unary \texttt{-}
  \item Function call, indexing, field access
\end{enumerate}

Use parentheses for clarity when precedence is not obvious:

\begin{lstlisting}
2 + 3 * 4       // 14  (multiplication first)
(2 + 3) * 4     // 20
\end{lstlisting}
